\ifdefined\gkver\else\def\gkver{student}\fi
\documentclass[\gkver]{gaokaozhenti}
\begin{document}

\chapter{2001年全国理综旧课程}
%% paperType: 理综物理部分

\begin{enumerate}
\item
%% number: 16
%% paperName: 2001年全国理综旧课程第 16 题 6 分
%% typeId: 1
%% type: 单选题
%% chapter: 光学
%% point: 光的干涉
%% method: 无
%% score: 6
%% degree: 650
%% duplicateId: 0
%% body:
市场上有种灯具俗称“冷光灯”，用它照射物品时能使被照物品处产生热效应大大降低，从而广泛地应用于博物馆、商店等处。这种灯降低热效应的原因之一是在灯泡后面放置的反光镜玻璃表面上镀一层薄膜（例如氟化镁），这种膜能消除不镀膜时玻璃表面反射回来的热效应最显著的红外线。以 $\lambda$ 表示此红外线的波长，则所镀薄膜的厚度最小应为\xzanswer{B}

\fourchoices[answer=B]
{$\frac{1}{8}\lambda$}
{$\frac{1}{4}\lambda$}
{$\frac{1}{2}\lambda$}
{$\lambda$}

%% answer: B
%% memo:
\memoanswer{
为减小反射热效应显著的红外线，则要求红外线在薄膜的前后表面反射后叠加作用减弱，即光程差为半波长的奇数倍，故膜的最小厚度为红外线在该膜中波长的 $\frac{1}{4}$，故 B 正确。
}

\item
%% number: 17
%% paperName: 2001年全国理综旧课程第 17 题 6 分
%% typeId: 1
%% type: 单选题
%% chapter: 气体性质
%% point: 气体图象
%% method: 无
%% score: 6
%% degree: 650
%% duplicateId: 0
%% body:
一定质量的理想气体由状态 A 经过图中所示过程变到状态 B。在此过程中气体的密度\xzanswer{A}

\onepicture[w=4cm]{\includesvg{figs/17}}

\fourchoices[answer=A]
{一直变小}
{一直变大}
{先变小后变大}
{先变大后变小}

%% answer: A
%% memo:
\memoanswer{
由图象可知，从 A 到 B，气体的温度变大，压强变小，由理想气体状态方程 $\frac{p_{A}V_{A}}{T_{A}}=\frac{p_{B}V_{B}}{T_{B}}$ 得，气体的体积变大，由 $\rho=\frac{M}{V}$ 得，气体的密度变小，故 A 正确。
}

\item
%% number: 18
%% paperName: 2001年全国理综旧课程第 18 题 6 分
%% typeId: 1
%% type: 单选题
%% chapter: 力物体的平衡
%% point: 物体系平衡
%% method: 无
%% score: 6
%% degree: 650
%% duplicateId: 0
%% body:
如图所示，在一粗糙水平面上有两个质量分别为 $m_{1}$ 和 $m_{2}$ 的木块 1 和 2，中间用一原长为 $l$、劲度系数为 $k$ 的轻弹簧连接起来，木块与地面间的滑动摩擦因数为 $\mu$。现用一水平力向右拉木块 2，当两木块一起匀速运动时两木块之间的距离是\xzanswer{A}

\onepicture[w=6cm]{figs/18.png}

\fourchoices[answer=A]
{$l+\frac{\mu}{k}m_{1}g$}
{$l+\frac{\mu}{k}(m_{1}+m_{2})g$}
{$l+\frac{\mu}{k}m_{2}g$}
{$l+\frac{\mu}{k}\cdot\frac{m_{1}m_{2}}{m_{1}+m_{2}}g$}

%% answer: A
%% memo:
\memoanswer{
\onepicture[align=right, w=8cm]{figs/18_an.png}

以连接体整体为研究对象，进行受力分析，其受力如图甲所示，对整体，由三力平衡得 $F-F_{1}-F_{2}=0$，其中，$F_{1}=\mu m_{1}g$，$F_{2}=\mu m_{2}g$；以木块 2 为研究对象，其受力如图乙所示，由三力平衡得 $F-F_{2}-F_{\text{弹}}=0$，其中，$F_{\text{弹}}$ 为弹簧的弹力，联立解得 $F_{\text{弹}}=\mu m_{1}g$；设弹簧的伸长量为 $x$，由胡克定律得 $F_{\text{弹}}=kx$，即 $x=\frac{\mu}{k}m_{1}g$，所以当两木块一起匀速运动时两木块之间的距离为 $l+x=l+\frac{\mu}{k}m_{1}g$，故 A 正确。
}

\item
%% number: 19
%% paperName: 2001年全国理综旧课程第 19 题 6 分
%% typeId: 1
%% type: 单选题
%% chapter: 直线运动
%% point: 运动的合成
%% method: 无
%% score: 6
%% degree: 450
%% duplicateId: 0
%% body:
在抗洪抢险中，战士驾驶摩托艇救人。假设江岸是平直的，洪水沿江向下游流去，水流速度为 $v_{1}$，摩托艇在静水中的航速为 $v_{2}$，战士救人的地点 A 离岸边最近处 O 的距离为 $d$。如战士想在最短时间内将人送上岸，则摩托艇登陆的地点离 O 点的距离为\xzanswer{C}

\fourchoices[answer=C]
{$\frac{dv_{2}}{\sqrt{v_{2}^{2}-v_{1}^{2}}}$}
{0}
{$\frac{dv_{1}}{v_{2}}$}
{$\frac{dv_{2}}{v_{1}}$}

%% answer: C
%% memo:
\memoanswer{
摩托艇要想在最短时间内到达对岸，其行驶方向要垂直于江岸，摩托艇的实际运动是相对于水的运动和随水流的运动的合运动，垂直于江岸方向的运动速度为 $v_{2}$，到达江岸所用的时间为 $t=\frac{d}{v_{2}}$；沿江岸方向的运动速度是水速 $v_{1}$，在相同的时间内，被水冲下的距离即为登陆点距离 O 点距离 $s=v_{1}t=\frac{dv_{1}}{v_{2}}$，故 C 正确。
}

\item
%% number: 20
%% paperName: 2001年全国理综旧课程第 20 题 6 分
%% typeId: 1
%% type: 单选题
%% chapter: 机械波
%% point: 波与振动图象
%% method: 无
%% score: 6
%% degree: 550
%% duplicateId: 0
%% body:
图 \ref{2001全国理综旧课程20a} 所示为一列简谐横波在 $t=20\Us$ 时的波形图，图 \ref{2001全国理综旧课程20b} 是这列波中 P 点的振动图线，那么该波的传播速度和传播方向是\xzanswer{B}

\twopicture[labelA=2001全国理综旧课程20a, labelB=2001全国理综旧课程20b, numA=1, numB=2, w=6cm]{\includesvg{figs/20a}}{\includesvg{figs/20b}}

\fourchoices[answer=B]
{$v=25\Ucms$，向左传播}
{$v=50\Ucms$，向左传播}
{$v=25\Ucms$，向右传播}
{$v=50\Ucms$，向右传播}

%% answer: B
%% memo:
\memoanswer{
由图 \ref{2001全国理综旧课程20a} 知波长为 $100\Ucm$，图 \ref{2001全国理综旧课程20b} 知波传播的周期为 $2\Us$，故 $v=\frac{\lambda}{T}=50\Ucms$。通过图 \ref{2001全国理综旧课程20b} 和振动周期性可知，在此时刻，P 点向上振动，所以横波向左传播，故 ACD 错误 B 正确。
}

\item
%% number: 21
%% paperName: 2001年全国理综旧课程第 21 题 6 分
%% typeId: 1
%% type: 单选题
%% chapter: 电场
%% point: 电势差与电场强度的关系
%% method: 无
%% score: 6
%% degree: 550
%% duplicateId: 0
%% body:
图中所示是一个平行板电容器，其电容为 $C$，带电量为 $Q$，上极板带正电。现将一个试探电荷 $q$ 由两极板间的 A 点移动到 B 点，如图所示。A、B 两点间的距离为 $s$，连线 AB 与极板间的夹角为 $30^\circ$，则电场力对试探电荷 $q$ 所做的功等于\xzanswer{C}

\onepicture[w=6cm]{\includesvg{figs/21}}

\fourchoices[answer=C]
{$\frac{qCs}{Qd}$}
{$\frac{qQs}{Cd}$}
{$\frac{qQs}{2Cd}$}
{$\frac{qCs}{2Qd}$}

%% answer: C
%% memo:
\memoanswer{
由题可知，电容器两平板间的电压 $U=\frac{Q}{C}$，平行板间的电场为 $E=\frac{U}{d}=\frac{Q}{Cd}$，电场对试探电荷 $q$ 做的功为 $W=Fs=qEs=\frac{qQ}{Cd}s\sin30^\circ$，整理得 $W=\frac{Qqs}{2Cd}$，故 C 正确。
}

\item
%% number: 22
%% paperName: 2001年全国理综旧课程第 22 题 6 分
%% typeId: 1
%% type: 单选题
%% chapter: 光学
%% point: 光的折射
%% method: 无
%% score: 6
%% degree: 650
%% duplicateId: 0
%% body:
如图所示，两块同样的玻璃直角三棱镜 ABC，两者的 AC 面是平行放置的，在它们之间是均匀的未知透明介质。一单色细光束 O 垂直于 AB 面入射，在图示的出射光线中\xzanswer{B}

\onepicture[w=6cm]{figs/22.png}

\fourchoices[answer=B]
{1、2、3（彼此平行）中的任一条都有可能}
{4、5、6（彼此平行）中的任一条都有可能}
{7、8、9（彼此平行）中的任一条都有可能}
{只能是 4、6 中的某一条}

%% answer: B
%% memo:
\memoanswer{
光线由左边三棱镜 AB 面射入棱镜，不改变方向；接着穿过两三棱镜间的未知透明介质进入右边的三棱镜，由于透明介质的两表面是平行的，因此它的光学特性相当于一块两面平行的玻璃砖，能使光线发生平行侧移，只是因为它两边的介质不是真空，而是折射率未知的玻璃，因此是否侧移以及侧移的方向无法确定（若未知介质的折射率 $n$ 与玻璃折射率 $n_{\text{玻}}$ 相等，不侧移；若 $n>n_{\text{玻}}$ 时，向上侧移；若 $n<n_{\text{玻}}$ 时，向下侧移），但至少可以确定方向没变，仍然与棱镜的 AB 面垂直，这样光线由右边三棱镜 AB 面射出棱镜时，不改变方向，应为 4、5、6 中的任意一条，故 B 正确。
}

\item
%% number: 23
%% paperName: 2001年全国理综旧课程第 23 题 6 分
%% typeId: 1
%% type: 单选题
%% chapter: 功和能
%% point: 功
%% method: 无
%% score: 6
%% degree: 450
%% duplicateId: 0
%% body:
下列一些说法正确的是\xzanswer{D}

①一质点受两个力作用且处于平衡状态（静止或匀速），这两个力在同一段时间内的冲量一定相同

②一质点受两个力作用且处于平衡状态（静止或匀速），这两个力在同一段时间内做的功或者都为零，或者大小相等符号相反

③在同样时间内，作用力和反作用力的功大小不一定相等，但正负号一定相反

④在同样时间内，作用力和反作用力的功大小不一定相等，正负号也不一定相反

\fourchoices[answer=D]
{①②}
{①③}
{②③}
{②④}

%% answer: D
%% memo:
\memoanswer{
因为处于平衡状态时，两个力大小相等方向相反，在同一段时间内冲量大小相等，但方向相反，说法①不正确；由恒力做功的知识可知，说法②正确；关于作用力和反作用力的功要认识到它们是作用在两个物体上，两个物体的位移可能不同，所以功可能不同，说法③不正确，说法④正确，故 D 正确。
}

\item
%% number: 24
%% paperName: 2001年全国理综旧课程第 24 题 6 分
%% typeId: 1
%% type: 单选题
%% chapter: 电磁感应
%% point: 电磁感应电路分析
%% method: 无
%% score: 6
%% degree: 250
%% duplicateId: 0
%% body:
电磁流量计广泛应用于测量可导电流体（如污水）在管中的流量（在单位时间内通过管内横截面的流体的体积）。为了简化，假设流量计是如图所示的横截面为长方形的一段管道，其中空部分的长、宽、高分别为图中的 $a$、$b$、$c$，流量计的两端与输送液体的管道相连接（图中虚线）。图中流量计的上下两面是金属材料，前后两面是绝缘材料，现于流量计所在处加磁感强度为 $B$ 的匀强磁场，磁场方向垂直于前后两面。当导电液体稳定地流经流量计时，在管外将流量计上、下两表面分别与一串接了电阻 $R$ 的电流表的两端连接，$I$ 表示测得的电流值。已知流体的电阻率为 $\rho$，不计电流表的内阻，则可求得流量为\xzanswer{A}

\onepicture[w=6cm]{\includesvg{figs/24}}

\fourchoices[answer=A]
{$\frac{I}{B}(bR+\rho\frac{c}{a})$}
{$\frac{I}{B}(aR+\rho\frac{b}{c})$}
{$\frac{I}{B}(cR+\rho\frac{a}{b})$}
{$\frac{I}{B}(R+\rho\frac{bc}{a})$}

%% answer: A
%% memo:
\memoanswer{
\onepicture[align=right, w=5cm]{figs/24_an.png}

流量计的工作原理如图所示，导电流体内的载流子在磁场的作用下向上下表面偏转，形成电势差 $U=Bvc$（其中 $v$ 为流体的速度），与外电路构成闭合电路时形成电流，其中流体内部存在一定的内阻 $r=\rho\frac{c}{ab}$，由闭合电路欧姆定律得 $U=I(R+r)$，解得流速 $v=\frac{I(R+\frac{\rho c}{ab})}{Bc}$，流量 $Q=bc\cdot v=\frac{I}{B}(bR+\rho\frac{c}{a})$，故 A 正确。
}

\item
%% number: 29
%% paperName: 2001年全国理综旧课程第 29 题 20 分
%% typeId: 4
%% type: 实验题
%% chapter: 电路
%% point: 电路实验
%% method: 无
%% score: 20
%% degree: 450
%% duplicateId: 0
%% body:
实验室中现有器材如实物图 \ref{2001全国理综旧课程29a} 所示，有：

电池 $E$，电动势约 $10\UV$，内阻约 $1\UO$；

电流表 A$_{1}$，量程 $10\UA$，内阻 $r_{1}$ 约为 $0.2\UO$；

电流表 A$_{2}$，量程 $300\UmA$，内阻 $r_{2}$ 约为 $5\UO$；

电流表 A$_{3}$，量程 $250\UmA$，内阻 $r_{3}$ 约为 $5\UO$；

电阻箱 $R_{1}$，最大阻值 $999.9\UO$，阻值最小改变量为 $0.1\UO$；

滑线变阻器 $R_{2}$，最大阻值 $100\UO$；开关 S；导线若干。

要求用图 \ref{2001全国理综旧课程29b} 所示的电路测定图中电流表 A 的内阻。

\begin{enumerate}
	\item
	在所给的三个电流表中，哪几个可用此电路精确测出其内阻？
	\item
	在可测的电流表中任选一个作为测量对象，在实物图上连成测量电路。
	\item
	你要读出的物理量是 \tkanswer[6cm]{A$_{2}$、A$_{3}$ 两电流表的读数 $I_{2}$、$I_{3}$ 和电阻箱 $R_{1}$ 的阻值 $R_{1}$}。

	用这些物理量表示待测内阻的计算公式是 \tkanswer[3cm]{$r_{3}=\frac{I_{2}-I_{3}}{I_{3}}R_{1}$}。
\end{enumerate}

\twopicture[numA=1, numB=2, labelA=2001全国理综旧课程29a, labelB=2001全国理综旧课程29b, wA=9cm, wB=5.5cm]{figs/29a.png}{figs/29b.png}

%% answer: 
\jdanswer{
\begin{enumerate}
	\item
	A$_{2}$、A$_{3}$
	\item
	若测 $r_{3}$，实物图如图：

	\onepicture[align=center, w=8cm]{figs/29_an.png}
	\item
	A$_{2}$、A$_{3}$ 两电流表的读数 $I_{2}$、$I_{3}$ 和电阻箱 $R_{1}$ 的阻值 $R_{1}$。

	$r_{3}=\frac{I_{2}-I_{3}}{I_{3}}R_{1}$
\end{enumerate}
}
%% memo:
\memoanswer{
若电阻箱 $R_{1}$ 和电流表 A 并联，通过的电流与阻值成反比，因此干路中的电流将大部分通过电流表 A，通过电流表 A$'$ 与电流表 A 的电流强度数值相差不大，两电流表的量程应相差不大，根据题设条件所提供的三只电流表只能选用 A$_{2}$、A$_{3}$ 电流表；另一方面整个电路的最大电阻约为 $110\UO$，最小电流约为 $91\UmA$，为了提高测量精确度，也应采用 A$_{2}$、A$_{3}$ 电流表。
}

\item
%% number: 30
%% paperName: 2001年全国理综旧课程第 30 题 14 分
%% typeId: 6
%% type: 计算题
%% chapter: 磁场
%% point: 带电粒子在磁场中的运动
%% method: 无
%% score: 14
%% degree: 550
%% duplicateId: 0
%% body:
如图是测量带电粒子质量的仪器工作原理示意图。设法是某有机化合物的气态分子导入图中所示的容器 A 中，使它受到电子束轰击，失去一个电子变成正一价的分子离子。分子离子从狭缝 $s_{1}$ 以很小的速度进入电压为 $U$ 的加速电场区（初速不计），加速后，再通过狭缝 $s_{2}$、$s_{3}$ 射入磁感强度为 $B$ 的匀强磁场，方向垂直于磁场区的界面 PQ。最后，分子离子打到感光片上，形成垂直于纸面而且平行于狭缝 $s_{3}$ 的细线。若测得细线到狭缝 $s_{3}$ 的距离为 $d$。导出分子离子的质量 $m$ 的表达式。

\onepicture[align=right, w=6cm]{figs/30.png}

%% answer: 
\jdanswer{
$m=\frac{eB^{2}d^{2}}{8U}$
}
%% memo:
\memoanswer{
\onepicture[align=right, w=5cm]{figs/30_an.png}

以 $m$、$q$ 表示离子的质量电量，以 $v$ 表示离子从狭缝 $s_{2}$ 射出时的速度，由功能关系可得

$eU=\frac{1}{2}mv^{2}$ ①

射入磁场后，在洛伦兹力作用下做圆周运动，由牛顿定律可得

$evB=m\frac{v^{2}}{R}$ ②

式中 $R$ 为圆的半径。感光片上的细黑线到 $s_{3}$ 缝的距离

$d=2R$ ③

解得

$m=\frac{eB^{2}d^{2}}{8U}$ ④

评分标准：本题 24 分，其中第（1）问 14 分，第（2）问 3 分，第（3）问 7 分。

第（1）问中，①、②、③式各 3 分，④式 5 分；
}

\item
%% number: 31
%% paperName: 2001年全国理综旧课程第 31 题 28 分
%% typeId: 6
%% type: 计算题
%% chapter: 曲线运动
%% point: 天体运动
%% method: 近似与估算
%% score: 28
%% degree: 350
%% duplicateId: 0
%% body:
太阳现正处于主序星演化阶段。它主要是由电子和 \ce{^1_1H}、\ce{^4_2He} 等原子核组成。维持太阳辐射的是它内部的核聚变反应，核反应方程是 \ce{2^{0}_{-1}e + 4^{1}_{1}H -> ^{4}_{2}He} + 释放的核能，这些核能最后转化为辐射能。根据目前关于恒星演化的理论，若由于聚变反应而使太阳中的 \ce{^1_1H} 核数目从现有数减少 $10\%$，太阳将离开主序星阶段而转入红巨星的演化阶段。为了简化，假定目前太阳全部由电子和 \ce{^1_1H} 核组成。

\begin{enumerate}
	\item
	为了研究太阳演化进程，需知道目前太阳的质量 $M$。已知地球半径 $R=6.4\times10^{6}\Um$，地球质量 $m=6.0\times10^{24}\Ukg$，日地中心的距离 $r=1.5\times10^{11}\Um$，地球表面处的重力加速度 $g=10\Umsq$，1 年约为 $3.2\times10^{7}\Us$，试估算目前太阳的质量 $M$。
	\item
	已知质子质量 $m_{\mathrm{p}}=1.6726\times10^{-27}\Ukg$，\ce{^4_2He} 质量 $m_{\alpha}=6.6458\times10^{-27}\Ukg$，电子质量 $m_{\mathrm{e}}=0.9\times10^{-30}\Ukg$，光速 $c=3\times10^{8}\Ums$。求每发生一次题中所述的核聚变反应所释放的核能。
	\item
	又知地球上与太阳垂直的每平方米截面上，每秒通过的太阳辐射能 $w=1.35\times10^{3}\UWmq$。试估算太阳继续保持在主序星阶段还有多少年的寿命。（估算结果只要求一位有效数字。）
\end{enumerate}

%% answer: 
\jdanswer{
\begin{enumerate}
	\item
	$M=2\times10^{30}\Ukg$
	\item
	$\Delta E=4.2\times10^{-12}\UJ$
	\item
	$t=1\times10^{10}$ 年 $=1$ 百亿年
\end{enumerate}
}
%% memo:
\memoanswer{
（1）估算太阳的质量 $M$。

设 $T$ 为地球绕日心运动的周期，则由万有引力定律和牛顿定律可知

$G\frac{mM}{r^{2}}=m\left(\frac{2\pi}{T}\right)^{2}r$ ①

地球表面处的重力加速度

$g=G\frac{m}{R^{2}}$ ②

由 ①、②式联立解得

$M=m\left(\frac{2\pi}{T}\right)^{2}\frac{r^{3}}{R^{2}g}$ ③

以题给数值代入，得 $M=2\times10^{30}\Ukg$ ④

（2）根据质量亏损和质能公式，该核反应每发生一次释放的核能为

$\Delta E=(4m_{\mathrm{p}}+2m_{\mathrm{e}}-m_{\alpha})c^{2}$ ⑤

代入数值，得

$\Delta E=4.2\times10^{-12}\UJ$ ⑥

（3）根据题给假设，在太阳继续保持在主序星阶段的时间内，发生题中所述的核聚变反应的次数为

$N=\frac{M}{4m_{\mathrm{p}}}\times10\%$ ⑦

因此，太阳总共辐射出的能量为

$E=N\cdot\Delta E$

设太阳辐射是各向同性的，则每秒内太阳向外放出的辐射能为

$P_{\text{太}}=4\pi r^{2}w$ ⑧

所以太阳继续保持在主序的时间为

$t=\frac{E}{P_{\text{太}}}$ ⑨

由以上各式解得

$t=\frac{0.1M(4m_{\mathrm{p}}+2m_{\mathrm{e}}-m_{\alpha})c^{2}}{4m_{\mathrm{p}}\times4\pi r^{2}w}$

以题给数据代入，并以年为单位，可得

$t=1\times10^{10}$ 年 $=1$ 百亿年 ⑩

评分标准：本题 28 分，其中第（1）问 14 分，第（2）问 7 分，第（3）问 7 分。

第（1）问中 ①、②两式各 3 分，③式 4 分，得出④式 4 分；

第（2）问中 ⑤式 4 分，⑥式 3 分；

第（3）问中 ⑦、⑧两式各 2 分，⑨式 2 分，⑩式 1 分。
}

\end{enumerate}

\end{document}
